% GENERATED FILE. Edit canonical lessons, then run npm run generate:books.
% canonical-source-hash: fnv1a64:8e307f6b16363440
% canonical-lessons: TE-C408-samajam, TE-C408-samajika, TE-C408-sangham, TE-C408-sabhyudu, TE-C408-vivaksa

\chapter{Society}
\label{ch:te-society}
\hlchaptermodality{408}

\begin{chapteropening}
\textbf{By the end of this chapter:} \emph{I can talk about society, associations and discrimination.}

Complete the last lesson of chapter 408: I can talk about society, associations and discrimination.
\end{chapteropening}

\section[\texorpdfstring{\te{సమాజం} (samājaṁ)}{సమాజం (samājaṁ)}]{\te{సమాజం} (samājaṁ) — society; a community}

\label{lesson:TE-C408-samajam}



\begin{quote}
Before the new one: say the Telugu for a claim, then the Telugu for a reform.
\end{quote}

\subsection*{You'll want to know: \te{సమాజం}}
\textbf{\te{సమాజం}} — \emph{samājaṁ} — \textquotedblleft{}society; a community\textquotedblright{}.

\begin{quote}
\te{సమాజంలో} \te{ప్రతి} \te{ఒక్కరికీ} \te{ఒక} \te{బాధ్యత} \te{ఉంది}.
\end{quote}

\begin{quote}
\emph{samājaṁlō prati okkarikī oka bādhyata undi.}
\end{quote}

\begin{quote}
Everyone has a responsibility in society.
\end{quote}

Also a smaller community: \te{మన} \te{సమాజం}, our community.

\begin{cousinweb}
\textbf{\te{సమాజం}} (\emph{samājaṁ}) is borrowed from Sanskrit \emph{samāja}, \textquotedblleft{}a festive gathering\textquotedblright{} (CDIAL 13207), which Monier-Williams also glosses \textquotedblleft{}an assembly, a society\textquotedblright{}.
\end{cousinweb}

\begin{grammarlens}[title={What you've built}]
One more word for giving reasons and opinions.
\end{grammarlens}

\subsection*{Your turn}
\begin{itemize}
\raggedright
  \item \emph{Say it:} \emph{samājaṁ}
  \item \emph{Say it:} \emph{samājaṁ}, once more
  \item \emph{Say it:} \emph{samājaṁlō prati okkarikī oka bādhyata undi.} — Everyone has a responsibility in society.
  \item \emph{Recall:} say the Telugu for a claim, then the Telugu for a reform, then say \emph{samājaṁ} again
\end{itemize}

\begin{tcolorbox}[breakable,colback=teal!4,colframe=teal!35!black,title={Before you move on}]
What does \te{సమాజం} mean? (\textquotedblleft{}Society; a community\textquotedblright{}.) Say it once more.
\end{tcolorbox}
\section[\texorpdfstring{\te{సామాజిక} (sāmājika)}{సామాజిక (sāmājika)}]{\te{సామాజిక} (sāmājika) — social}

\label{lesson:TE-C408-samajika}



\begin{quote}
Before the new one: say the Telugu for a reform, then the Telugu for society.
\end{quote}

\subsection*{You'll want to know: \te{సామాజిక}}
\textbf{\te{సామాజిక}} — \emph{sāmājika} — \textquotedblleft{}social\textquotedblright{}.

\begin{quote}
\te{పేదరికం} \te{ఒక} \te{పెద్ద} \te{సామాజిక} \te{సమస్య}.
\end{quote}

\begin{quote}
\emph{pēdarikaṁ oka pedda sāmājika samasya.}
\end{quote}

\begin{quote}
Poverty is a big social problem.
\end{quote}

It stands before a noun: \te{సామాజిక} \te{సమస్య}, a social problem.

\begin{cousinweb}
\textbf{\te{సామాజిక}} (\emph{sāmājika}) is borrowed from Sanskrit \emph{sāmājika}, \textquotedblleft{}relating to an assembly\textquotedblright{} (Monier-Williams), from \emph{samāja}, the word Telugu has as \te{సమాజం}.
\end{cousinweb}

\begin{grammarlens}[title={What you've built}]
One more word for giving reasons and opinions.
\end{grammarlens}

\subsection*{Your turn}
\begin{itemize}
\raggedright
  \item \emph{Say it:} \emph{sāmājika}
  \item \emph{Say it:} \emph{sāmājika}, once more
  \item \emph{Say it:} \emph{pēdarikaṁ oka pedda sāmājika samasya.} — Poverty is a big social problem.
  \item \emph{Recall:} say the Telugu for a reform, then the Telugu for society, then say \emph{sāmājika} again
\end{itemize}

\begin{tcolorbox}[breakable,colback=teal!4,colframe=teal!35!black,title={Before you move on}]
What does \te{సామాజిక} mean? (\textquotedblleft{}Social\textquotedblright{}.) Say it once more.
\end{tcolorbox}
\section[\texorpdfstring{\te{సంఘం} (saṅghaṁ)}{సంఘం (saṅghaṁ)}]{\te{సంఘం} (saṅghaṁ) — an association; a union}

\label{lesson:TE-C408-sangham}



\begin{quote}
Before the new one: say the Telugu for society, then the Telugu for social.
\end{quote}

\subsection*{You'll want to know: \te{సంఘం}}
\textbf{\te{సంఘం}} — \emph{saṅghaṁ} — \textquotedblleft{}an association; a union\textquotedblright{}.

\begin{quote}
\te{మా} \te{వీధిలో} \te{ఒక} \te{యువజన} \te{సంఘం} \te{ఉంది}.
\end{quote}

\begin{quote}
\emph{mā vīdhilō oka yuvajana saṅghaṁ undi.}
\end{quote}

\begin{quote}
There is a youth association in our street.
\end{quote}

Plural: \te{సంఘాలు}, \emph{saṅghālu}.

Also a trade union. Plural \te{సంఘాలు}.

\begin{cousinweb}
\textbf{\te{సంఘం}} (\emph{saṅghaṁ}) is borrowed from Sanskrit \emph{saṁgha}, \textquotedblleft{}an association, a community\textquotedblright{} (CDIAL 12854).
\end{cousinweb}

\begin{grammarlens}[title={What you've built}]
One more word for giving reasons and opinions.
\end{grammarlens}

\subsection*{Your turn}
\begin{itemize}
\raggedright
  \item \emph{Say it:} \emph{saṅghaṁ}
  \item \emph{Say it:} \emph{saṅghaṁ}, once more
  \item \emph{Say it:} \emph{mā vīdhilō oka yuvajana saṅghaṁ undi.} — There is a youth association in our street.
  \item \emph{Recall:} say the Telugu for society, then the Telugu for social, then say \emph{saṅghaṁ} again
\end{itemize}

\begin{tcolorbox}[breakable,colback=teal!4,colframe=teal!35!black,title={Before you move on}]
What does \te{సంఘం} mean? (\textquotedblleft{}An association; a union\textquotedblright{}.) Say it once more.
\end{tcolorbox}
\section[\texorpdfstring{\te{సభ్యుడు} (sabhyuḍu)}{సభ్యుడు (sabhyuḍu)}]{\te{సభ్యుడు} (sabhyuḍu) — a member}

\label{lesson:TE-C408-sabhyudu}



\begin{quote}
Before the new one: say the Telugu for social, then the Telugu for an association.
\end{quote}

\subsection*{You'll want to know: \te{సభ్యుడు}}
\textbf{\te{సభ్యుడు}} — \emph{sabhyuḍu} — \textquotedblleft{}a member\textquotedblright{}.

\begin{quote}
\te{నేను} \te{ఆ} \te{సంఘంలో} \te{సభ్యుడిని}.
\end{quote}

\begin{quote}
\emph{nēnu ā saṅghaṁlō sabhyuḍini.}
\end{quote}

\begin{quote}
I am a member of that association.
\end{quote}

Plural: \te{సభ్యులు}, \emph{sabhyulu}.

A woman member is \te{సభ్యురాలు}.

\begin{cousinweb}
\textbf{\te{సభ్యుడు}} (\emph{sabhyuḍu}) is borrowed from Sanskrit \emph{sabhya}, \textquotedblleft{}being in an assembly\textquotedblright{} (CDIAL 13171), from \emph{sabhā}, \textquotedblleft{}an assembly\textquotedblright{}, with the Telugu masculine ending -\te{డు}.
\end{cousinweb}

\begin{grammarlens}[title={What you've built}]
One more word for giving reasons and opinions.
\end{grammarlens}

\subsection*{Your turn}
\begin{itemize}
\raggedright
  \item \emph{Say it:} \emph{sabhyuḍu}
  \item \emph{Say it:} \emph{sabhyuḍu}, once more
  \item \emph{Say it:} \emph{nēnu ā saṅghaṁlō sabhyuḍini.} — I am a member of that association.
  \item \emph{Recall:} say the Telugu for social, then the Telugu for an association, then say \emph{sabhyuḍu} again
\end{itemize}

\begin{tcolorbox}[breakable,colback=teal!4,colframe=teal!35!black,title={Before you move on}]
What does \te{సభ్యుడు} mean? (\textquotedblleft{}A member\textquotedblright{}.) Say it once more.
\end{tcolorbox}
\section[\texorpdfstring{\te{వివక్ష} (vivakṣa)}{వివక్ష (vivakṣa)}]{\te{వివక్ష} (vivakṣa) — discrimination}

\label{lesson:TE-C408-vivaksa}



\begin{quote}
Before the new one: say the Telugu for an association, then the Telugu for a member.
\end{quote}

\subsection*{You'll want to know: \te{వివక్ష}}
\textbf{\te{వివక్ష}} — \emph{vivakṣa} — \textquotedblleft{}discrimination\textquotedblright{}.

\begin{quote}
\te{ఆడపిల్లల} \te{మీద} \te{వివక్ష} \te{తప్పు}.
\end{quote}

\begin{quote}
\emph{āḍapillala mīda vivakṣa tappu.}
\end{quote}

\begin{quote}
Discrimination against girls is wrong.
\end{quote}

The group treated unfairly takes \te{మీద} or \te{పట్ల}.

\begin{cousinweb}
\textbf{\te{వివక్ష}} (\emph{vivakṣa}) is borrowed from Sanskrit \emph{vivakṣā}, \textquotedblleft{}the wish to speak; meaning, intention\textquotedblright{} (Monier-Williams); the sense \textquotedblleft{}discrimination\textquotedblright{} is the modern Telugu one.
\end{cousinweb}

\begin{grammarlens}[title={What you've built}]
One more word for giving reasons and opinions.
\end{grammarlens}

\subsection*{Your turn}
\begin{itemize}
\raggedright
  \item \emph{Say it:} \emph{vivakṣa}
  \item \emph{Say it:} \emph{vivakṣa}, once more
  \item \emph{Say it:} \emph{āḍapillala mīda vivakṣa tappu.} — Discrimination against girls is wrong.
  \item \emph{Recall:} say the Telugu for an association, then the Telugu for a member, then say \emph{vivakṣa} again
\end{itemize}

\begin{tcolorbox}[breakable,colback=teal!4,colframe=teal!35!black,title={Before you move on}]
What does \te{వివక్ష} mean? (\textquotedblleft{}Discrimination\textquotedblright{}.) Say it once more.
\end{tcolorbox}
